%% Chapter 5 — 2026 Current Assessment (T₁)
\chapter{2026 Current Assessment (\Tone)}
\label{ch:current}

\section{Introduction}

This chapter presents the habitat classification and biodiversity unit calculation for the 2026 present-day epoch (\Tone). The assessment represents conditions during the growing season of 2025--2026, captured by Sentinel-2A/B and Landsat~8/9 imagery composites. The identical classification pipeline and BM\,4.0 scoring methodology described in Chapter~\ref{ch:methodology} was applied to ensure direct comparability with the \Tzero{} baseline.

\begin{confidencebox}{HIGH}
Both the 2016 and 2026 classifications use the same training label set (updated for known changes only), the same Random Forest hyperparameters, and the same spectral feature set. This design minimises systematic bias in the delta calculation.
\end{confidencebox}


\section{Imagery Composite Quality}

The 2026 growing-season composite was generated from \dataplaceholder{N SCENES S2} Sentinel-2 scenes (both 2A and 2B, providing 5-day revisit) and \dataplaceholder{N SCENES L8} Landsat~8/9 scenes acquired between 1~May and 30~September 2025. The median composite achieves wall-to-wall coverage with an effective cloud-free observation density of \dataplaceholder{N OBS} observations per pixel.

\begin{figure}[H]
\centering
\begin{tikzpicture}
  \draw[dreamlab-light,fill=dreamlab-light,rounded corners] (0,0) rectangle (12,7);
  \node[dreamlab-grey,font=\sffamily\large] at (6,3.5) {Figure placeholder: 2026 true-colour composite};
  \node[dreamlab-grey,font=\sffamily\footnotesize] at (6,2.5) {../output/figures/composite\_2026\_rgb.png};
\end{tikzpicture}
\caption{Sentinel-2 true-colour composite (B4-B3-B2) for the 2025/2026 growing season, covering the 1\,km study area.}
\label{fig:composite-2026}
\end{figure}

\subsection{Comparison with 2016 Composite}

Visual comparison of the two composites reveals several notable changes:

\begin{itemize}[nosep]
  \item \dataplaceholder{CHANGE 1 — describe visible change in reflectance pattern};
  \item \dataplaceholder{CHANGE 2 — describe change in vegetation extent};
  \item \dataplaceholder{CHANGE 3 — describe change in field boundary patterns};
  \item \dataplaceholder{CHANGE 4 — describe any new or removed features}.
\end{itemize}

\begin{figure}[H]
\centering
\begin{tikzpicture}
  \draw[dreamlab-light,fill=dreamlab-light,rounded corners] (0,0) rectangle (12,7);
  \node[dreamlab-grey,font=\sffamily\large] at (6,3.5) {Figure placeholder: Side-by-side 2016 vs 2026 comparison};
  \node[dreamlab-grey,font=\sffamily\footnotesize] at (6,2.5) {../output/figures/composite\_comparison\_2016\_2026.png};
\end{tikzpicture}
\caption{Side-by-side comparison of growing-season true-colour composites: 2016 (left) and 2026 (right).}
\label{fig:composite-comparison}
\end{figure}


\section{Parcel Framework}

The 2026 parcel framework contains \dataplaceholder{N PARCELS} habitat parcels. Where parcel boundaries have changed between the OS MasterMap 2016 and 2026 snapshots (due to subdivision, merger, or boundary revision), a harmonisation procedure was applied:

\begin{enumerate}[nosep]
  \item Parcels present in both epochs with unchanged boundaries are directly compared;
  \item Subdivided parcels are tracked via geometric intersection, with the parent parcel's \Tzero{} BU apportioned by area;
  \item Merged parcels combine the constituent \Tzero{} BU values;
  \item New parcels (not present in 2016 MasterMap) are assigned a \Tzero{} BU of zero.
\end{enumerate}

\begin{table}[H]
\centering
\caption{Parcel framework harmonisation summary.}
\label{tab:parcel-harmonisation}
\begin{tabular}{lr}
\toprule
\textbf{Category} & \textbf{Count} \\
\midrule
Unchanged parcels       & \dataplaceholder{N} \\
Subdivided parcels      & \dataplaceholder{N} \\
Merged parcels          & \dataplaceholder{N} \\
New parcels (2026 only) & \dataplaceholder{N} \\
Removed (built over)    & \dataplaceholder{N} \\
\midrule
\textbf{Total T\textsubscript{1} parcels} & \dataplaceholder{TOTAL} \\
\bottomrule
\end{tabular}
\end{table}


\section{Habitat Classification Results}

\subsection{Classification Map}

\begin{figure}[H]
\centering
\begin{tikzpicture}
  \draw[dreamlab-light,fill=dreamlab-light,rounded corners] (0,0) rectangle (12,8);
  \node[dreamlab-grey,font=\sffamily\large] at (6,4) {Figure placeholder: 2026 habitat classification map};
  \node[dreamlab-grey,font=\sffamily\footnotesize] at (6,3) {../output/figures/habitat\_map\_2026.png};
\end{tikzpicture}
\caption{UK Habitat Classification map for the 2026 present-day epoch (\Tone), showing Level~3 habitat assignments per parcel. Colour coding follows the UKHab standard palette.}
\label{fig:habitat-map-2026}
\end{figure}

\subsection{Habitat Composition}

\begin{longtable}{p{50mm}rrp{30mm}}
\caption{Habitat composition at \Tone{} (2026) — area-based summary by UKHab Level~3 class.}
\label{tab:habitat-composition-2026} \\
\toprule
\textbf{UKHab Level 3 Class} & \textbf{Area (ha)} & \textbf{\% of Total} & \textbf{BM\,4.0 Distinctiveness} \\
\midrule
\endfirsthead
\toprule
\textbf{UKHab Level 3 Class} & \textbf{Area (ha)} & \textbf{\% of Total} & \textbf{BM\,4.0 Distinctiveness} \\
\midrule
\endhead
\midrule
\multicolumn{4}{r}{\small\itshape Continued on next page\ldots} \\
\bottomrule
\endfoot
\bottomrule
\endlastfoot

g3c --- Other neutral grassland         & \dataplaceholder{AREA} & \dataplaceholder{PCT} & Medium (4) \\
g4 --- Modified grassland               & \dataplaceholder{AREA} & \dataplaceholder{PCT} & Low (2) \\
g1a --- Lowland dry acid grassland       & \dataplaceholder{AREA} & \dataplaceholder{PCT} & High (6) \\
g3a --- Lowland calcareous grassland     & \dataplaceholder{AREA} & \dataplaceholder{PCT} & V.~High (8) \\
g6 --- Bracken                           & \dataplaceholder{AREA} & \dataplaceholder{PCT} & Low (2) \\
w1f --- Lowland mixed deciduous woodland & \dataplaceholder{AREA} & \dataplaceholder{PCT} & High (6) \\
w1g --- Other broadleaved woodland       & \dataplaceholder{AREA} & \dataplaceholder{PCT} & Medium (4) \\
w2c --- Other Scot's pine woodland       & \dataplaceholder{AREA} & \dataplaceholder{PCT} & Medium (4) \\
h1a --- Dwarf shrub heath (dry)          & \dataplaceholder{AREA} & \dataplaceholder{PCT} & High (6) \\
c1a --- Arable field margins             & \dataplaceholder{AREA} & \dataplaceholder{PCT} & Medium (4) \\
c1 --- Arable (cereal crops)             & \dataplaceholder{AREA} & \dataplaceholder{PCT} & Low (2) \\
r1 --- Standing open water               & \dataplaceholder{AREA} & \dataplaceholder{PCT} & Medium (4) \\
u1b --- Developed land; sealed surface   & \dataplaceholder{AREA} & \dataplaceholder{PCT} & V.~Low (0) \\
u1c --- Artificial unvegetated           & \dataplaceholder{AREA} & \dataplaceholder{PCT} & V.~Low (0) \\
\midrule
\textbf{Total}                           & \dataplaceholder{TOTAL} & 100.0 & --- \\
\end{longtable}

\subsection{Habitat Transitions Since 2016}

The following table summarises the most significant habitat class transitions between \Tzero{} and \Tone:

\begin{table}[H]
\centering
\caption{Principal habitat class transitions between 2016 and 2026.}
\label{tab:habitat-transitions}
\begin{tabular}{p{35mm}p{35mm}rp{35mm}}
\toprule
\textbf{From (2016)} & \textbf{To (2026)} & \textbf{Area (ha)} & \textbf{Interpretation} \\
\midrule
g4 (Modified grassland) & g3c (Neutral grassland) & \dataplaceholder{A} & Grassland improvement through reduced management intensity \\
c1 (Arable) & g4 (Modified grassland) & \dataplaceholder{A} & Arable reversion to grassland \\
g3c (Neutral grassland) & w1g (Broadleaved woodland) & \dataplaceholder{A} & Natural succession / woodland encroachment \\
g1a (Acid grassland) & g6 (Bracken) & \dataplaceholder{A} & Bracken encroachment on acid grassland \\
g4 (Modified grassland) & u1b (Developed) & \dataplaceholder{A} & Development / land conversion \\
\bottomrule
\end{tabular}
\end{table}


\section{Condition Assessment}

\subsection{Condition Distribution}

\begin{table}[H]
\centering
\caption{Distribution of condition scores at \Tone{} (2026) across all habitat parcels.}
\label{tab:condition-2026}
\begin{tabular}{lrrrr}
\toprule
\textbf{Condition} & \textbf{N Parcels} & \textbf{\%} & \textbf{Area (ha)} & \textbf{\% Area} \\
\midrule
Good        & \dataplaceholder{N} & \dataplaceholder{PCT} & \dataplaceholder{AREA} & \dataplaceholder{PCT} \\
Fairly Good & \dataplaceholder{N} & \dataplaceholder{PCT} & \dataplaceholder{AREA} & \dataplaceholder{PCT} \\
Moderate    & \dataplaceholder{N} & \dataplaceholder{PCT} & \dataplaceholder{AREA} & \dataplaceholder{PCT} \\
Fairly Poor & \dataplaceholder{N} & \dataplaceholder{PCT} & \dataplaceholder{AREA} & \dataplaceholder{PCT} \\
Poor        & \dataplaceholder{N} & \dataplaceholder{PCT} & \dataplaceholder{AREA} & \dataplaceholder{PCT} \\
N/A (built) & \dataplaceholder{N} & \dataplaceholder{PCT} & \dataplaceholder{AREA} & \dataplaceholder{PCT} \\
\midrule
\textbf{Total} & \dataplaceholder{TOTAL} & 100.0 & \dataplaceholder{TOTAL} & 100.0 \\
\bottomrule
\end{tabular}
\end{table}

\begin{figure}[H]
\centering
\begin{tikzpicture}
  \draw[dreamlab-light,fill=dreamlab-light,rounded corners] (0,0) rectangle (12,7);
  \node[dreamlab-grey,font=\sffamily\large] at (6,3.5) {Figure placeholder: 2026 condition map};
  \node[dreamlab-grey,font=\sffamily\footnotesize] at (6,2.5) {../output/figures/condition\_map\_2026.png};
\end{tikzpicture}
\caption{Parcel-level condition assessment at \Tone{} (2026).}
\label{fig:condition-map-2026}
\end{figure}

\subsection{Condition Change Summary}

\begin{table}[H]
\centering
\caption{Condition change between \Tzero{} and \Tone{} for parcels with stable habitat classification.}
\label{tab:condition-change}
\begin{tabular}{lrr}
\toprule
\textbf{Condition Trajectory} & \textbf{N Parcels} & \textbf{Area (ha)} \\
\midrule
Improved ($\geq$1 category)    & \dataplaceholder{N} & \dataplaceholder{AREA} \\
Stable (no change)             & \dataplaceholder{N} & \dataplaceholder{AREA} \\
Declined ($\geq$1 category)   & \dataplaceholder{N} & \dataplaceholder{AREA} \\
\bottomrule
\end{tabular}
\end{table}


\section{Strategic Significance}

\begin{table}[H]
\centering
\caption{Strategic significance distribution at \Tone{} (2026).}
\label{tab:strategic-2026}
\begin{tabular}{lrrr}
\toprule
\textbf{Strategic Significance} & \textbf{N Parcels} & \textbf{Area (ha)} & \textbf{Multiplier} \\
\midrule
High   & \dataplaceholder{N} & \dataplaceholder{AREA} & 1.15 \\
Medium & \dataplaceholder{N} & \dataplaceholder{AREA} & 1.10 \\
Low    & \dataplaceholder{N} & \dataplaceholder{AREA} & 1.00 \\
\midrule
\textbf{Total} & \dataplaceholder{TOTAL} & \dataplaceholder{TOTAL} & --- \\
\bottomrule
\end{tabular}
\end{table}


\section{Biodiversity Unit Calculation}

\subsection{BU Summary by Habitat Type}

\begin{longtable}{p{50mm}rrrrr}
\caption{Biodiversity unit summary by habitat type at \Tone{} (2026).}
\label{tab:bu-by-habitat-2026} \\
\toprule
\textbf{Habitat} & \textbf{Area} & \textbf{Distinct.} & \textbf{Cond.} & \textbf{Strat.} & \textbf{BU} \\
 & \textbf{(ha)} & & \textbf{(avg)} & \textbf{(avg)} & \\
\midrule
\endfirsthead
\toprule
\textbf{Habitat} & \textbf{Area} & \textbf{Distinct.} & \textbf{Cond.} & \textbf{Strat.} & \textbf{BU} \\
 & \textbf{(ha)} & & \textbf{(avg)} & \textbf{(avg)} & \\
\midrule
\endhead
\bottomrule
\endlastfoot

g3c — Other neutral grassland        & \dataplaceholder{A} & 4 & \dataplaceholder{C} & \dataplaceholder{S} & \dataplaceholder{BU} \\
g4 — Modified grassland              & \dataplaceholder{A} & 2 & \dataplaceholder{C} & \dataplaceholder{S} & \dataplaceholder{BU} \\
g1a — Lowland dry acid grassland     & \dataplaceholder{A} & 6 & \dataplaceholder{C} & \dataplaceholder{S} & \dataplaceholder{BU} \\
g3a — Lowland calcareous grassland   & \dataplaceholder{A} & 8 & \dataplaceholder{C} & \dataplaceholder{S} & \dataplaceholder{BU} \\
w1f — Lowland mixed deciduous wdl    & \dataplaceholder{A} & 6 & \dataplaceholder{C} & \dataplaceholder{S} & \dataplaceholder{BU} \\
w1g — Other broadleaved woodland     & \dataplaceholder{A} & 4 & \dataplaceholder{C} & \dataplaceholder{S} & \dataplaceholder{BU} \\
h1a — Dwarf shrub heath (dry)        & \dataplaceholder{A} & 6 & \dataplaceholder{C} & \dataplaceholder{S} & \dataplaceholder{BU} \\
c1 — Arable                          & \dataplaceholder{A} & 2 & \dataplaceholder{C} & \dataplaceholder{S} & \dataplaceholder{BU} \\
r1 — Standing open water             & \dataplaceholder{A} & 4 & \dataplaceholder{C} & \dataplaceholder{S} & \dataplaceholder{BU} \\
u1b — Developed; sealed surface      & \dataplaceholder{A} & 0 & --- & \dataplaceholder{S} & 0.00 \\
\midrule
\textbf{Total}                        & \dataplaceholder{TOTAL} & --- & --- & --- & \dataplaceholder{TOTAL BU} \\
\end{longtable}

\subsection{Spatial Distribution of Biodiversity Value}

\begin{figure}[H]
\centering
\begin{tikzpicture}
  \draw[dreamlab-light,fill=dreamlab-light,rounded corners] (0,0) rectangle (12,8);
  \node[dreamlab-grey,font=\sffamily\large] at (6,4) {Figure placeholder: 2026 BU choropleth map};
  \node[dreamlab-grey,font=\sffamily\footnotesize] at (6,3) {../output/figures/bu\_choropleth\_2026.png};
\end{tikzpicture}
\caption{Parcel-level biodiversity unit values at \Tone{} (2026), displayed as a graduated colour map.}
\label{fig:bu-choropleth-2026}
\end{figure}

\begin{confidencebox}{HIGH}
The total present-day biodiversity value of \dataplaceholder{TOTAL BU} area-based biodiversity units represents a \dataplaceholder{DIRECTION} of \dataplaceholder{ABS DELTA} BU from the 2016 baseline. The full delta analysis with uncertainty bounds is presented in Chapter~\ref{ch:delta}.
\end{confidencebox}
